% ===== Block E: PDF p.21–26 | 原文 §2.3（后半：$\psi$ 类相交数、Witten 猜想、Kontsevich 定理、拓扑递归与矩阵模型） | 唯一编号公式 2.62–2.79 =====
% 本块接 Block D 的式 (2.61) 之后开始；(2.62) 为首条编号公式，(2.79) 为末条。
% 编号环境：习题 2.9（p21）、习题 2.10（p22）、定理 2.6（p22）、习题 2.11（p22–p23）。
% 图：Figure 13（p24）、Figure 14（p25）从略；表：Table 1（p23）照抄重排。
% 注意：开头续接 Block D 中习题 2.8 的编号列表（原文 (1)(2)(3)(4)），故用 enumerate 续编号。

% [原文 p.21]
\begin{enumerate}[label=(\arabic*),start=2]
\item 对任意线丛 $\mathcal{L}$，有 $\mathrm{c}_1(\mathcal{L})=[Z-P]$，其中 $Z$ 与 $P$ 分别是 $\mathcal{L}$ 的某个一般亚纯截面的零点除子与极点除子，而 $[\cdot]$ 表示 Poincaré 对偶\footnote{轨形（orbifold）的 Poincaré 对偶涉及自同构群。更确切地说，若 $Z$ 是 $X$ 的子轨形、其底拓扑空间为 $\hat Z$，则 $[Z]=\frac{1}{|G|}[\hat Z]$，其中 $G$ 是 $\hat Z$ 中一般点的自同构群。}。
% [Tutor Note: 原文此处是全文第 2 号脚注（脚注计数器与正文分开），译注本改用 \footnote 自动编号。这个 1/|G| 因子不是装饰：在 $\Mc_{1,1}$ 上一般点的自同构群为 $\Z/2$（椭圆曲线的 $-1$ 自同构），而 Lagrange 关系给出 $12\lambda_1=\delta_1$、$\int_{\Mc_{1,1}}\delta_1=\frac{1}{2}$（边界点也是 $\Z/2$-轨形点），于是 $\int_{\Mc_{1,1}}\psi_1=\int\lambda_1=\frac{1}{2}\cdot\frac{1}{12}=\frac{1}{24}$；若漏掉该 1/2 就会得到 1/12。这正是之后 (2.75) 与 $\langle\tau_1\rangle_1=\frac{1}{24}$ 的来源。]
\item 考虑余切线丛 $\mathcal{L}_1^{\otimes k}\to\Mc_{1,1}$。$\mathcal{L}_1^{\otimes k}$ 的全纯截面构成的向量空间与权 $k$ 模形式构成的向量空间之间存在一个典范等同。
% [Tutor Note: 这一条与 (2.62) 直接呼应：$\mathcal{L}_1^{\otimes k}$ 的全纯截面就是权 $k$ 的模形式，于是 (2.62) 中的 $G_4,G_6,\tilde G_{12}$ 分别是 $\mathcal{L}_1^{\otimes 4},\mathcal{L}_1^{\otimes 6},\mathcal{L}_1^{\otimes 12}$ 的截面；$\tilde G_{12}$ 只在 $i\infty$ 处有一个单零点，据此 $\deg\mathcal{L}_1=\frac{1}{12}\cdot(\tilde G_{12}$ 在 $\ii\infty$ 处按轨形权计的零点阶数$)=\frac{1}{24}$，与上一条批注的算法一致。]
\item 下列（若干）Eisenstein 级数的组合
\begin{align}
G_4(\tau) &= \sum_{\lambda\in(\Z+\tau\Z)\setminus\{0\}}\frac{1}{\lambda^4}\,, \nonumber\\
G_6(\tau) &= \sum_{\lambda\in(\Z+\tau\Z)\setminus\{0\}}\frac{1}{\lambda^6}\,, \tag{2.62}\\
\tilde{G}_{12}(\tau) &= \left(\frac{G_4(\tau)}{2\zeta(4)}\right)^{3}-\left(\frac{G_6(\tau)}{2\zeta(6)}\right)^{2}, \nonumber
\end{align}
分别是权 $4$、$6$、$12$ 的模形式。此外，它们分别在 $\tau=\frac{1+\ii\sqrt3}{2}$、$\tau=\ii$ 与 $\tau=+\ii\infty$ 处有唯一的单零点。
% [Tutor Note: 用标准归一化核对 (2.62)：$G_{2k}=2\zeta(2k)E_{2k}$（$E_{2k}$ 为归一化 Eisenstein 级数），故 $G_4/(2\zeta(4))=E_4$、$G_6/(2\zeta(6))=E_6$，于是 $\tilde G_{12}=E_4^3-E_6^2=1728\Delta$。而 $E_4(\rho)=0$（$\rho=\ee^{\ii\pi/3}=\frac{1+\ii\sqrt3}{2}$）、$E_6(\ii)=0$、$\Delta$ 在基本域内只在尖点 $\ii\infty$ 处有单零点——三个零点位置与原文的断言逐个吻合。又 $\zeta(4)=\frac{\pi^4}{90}$、$\zeta(6)=\frac{\pi^6}{945}$，可用来检查这两个 $2\zeta$ 分母的归一化。]
% [Tutor Note: 原文把三行写在同一编号下、编号印在中间一行（$G_6$ 那行）右侧；译文用 align 复现此排法，第一、三行用 \nonumber。]
\end{enumerate}

我们现在可以陈述 Witten 猜想。首先，让我们把 Witten 的关联函数打包成单个生成级数：设 $t_d$（$d\ge0$）是一组形式变量，并令
\begin{equation}
Z(t_0,t_1,t_2,\ldots;\hbar)=\exp\left(\sum_{\substack{g\ge0,\,n\ge1\\ 2g-2+n>0}}\frac{\hbar^{2g-2+n}}{n!}\sum_{d_1,\ldots,d_n\ge0}\langle\tau_{d_1}\cdots\tau_{d_n}\rangle_g\prod_{i=1}^{n}t_{d_i}\right). \tag{2.63}
\end{equation}
% [Tutor Note: 求和条件 $2g-2+n>0$ 恰好排除 $(g,n)=(0,1)$ 与 $(0,2)$ 这两种 $\sum_i d_i=3g-3+n$ 无解（或需理解为负次幂）的情形；表 1 中列出的所有 $(g,n)$ 都满足它。最低阶项可用 $\langle\tau_0^3\rangle_0=1$ 写出：$Z=1+\hbar\frac{t_0^3}{6}+\mathrm{O}(\hbar^2)$，下面核对 $L_{-1}Z=0$ 与 $L_0Z=0$ 时都要用到这一项。]
生成级数 $Z$ 作为拓扑二维量子引力中的配分函数出现。弦方程与膨胀（dilaton）方程可以写成如下微分算子湮灭 $Z$ 的形式。

\begin{exercise}[2.9]
定义微分算子
\begin{equation}
L_{-1}=\hbar\frac{\partial}{\partial t_0}-\hbar^2\left(\sum_{k\ge1}t_k\frac{\partial}{\partial t_{k-1}}+\frac{t_0^2}{2}\right), \tag{2.64}
\end{equation}
\begin{equation}
L_0=\hbar\frac{\partial}{\partial t_1}-\hbar^2\left(\sum_{k\ge0}\frac{2k+1}{3}t_k\frac{\partial}{\partial t_k}+\frac{1}{24}\right). \tag{2.65}
\end{equation}
证明以下结论：
\begin{itemize}
\item 弦方程与 $\langle\tau_0^3\rangle_0=1$ 等价于方程 $L_{-1}Z=0$。
\item 膨胀方程与 $\langle\tau_1\rangle_1=\frac{1}{24}$ 等价于方程 $L_0Z=0$。
\end{itemize}
\end{exercise}
% [Tutor Note: (2.65) 中的系数 $1/3$ 可直接核对：在 $\langle\tau_{d_1}\cdots\tau_{d_n}\rangle_g\ne0$ 的单项式上（此时 $\sum_i d_i=3g-3+n$），$\sum_k\frac{2k+1}{3}t_k\partial_k$ 给出 $\frac{1}{3}\sum_i(2d_i+1)=\frac{1}{3}\bigl(2(3g-3+n)+n\bigr)=2g-2+n$，正是膨胀方程右端的因子；把 3 换成任何别的数都与膨胀方程不符。]
% [Tutor Note: (2.65) 中的 $1/24$ 也不是自由参数。写 $Z=\ee^F$，$Z$ 中含 $\hbar\,t_1\langle\tau_1\rangle_1$ 项；$L_0Z=0$ 的 $\hbar^2$ 阶、与 $t$ 无关的部分给出 $\hbar^2\langle\tau_1\rangle_1-\hbar^2\frac{1}{24}=0$，即 $\langle\tau_1\rangle_1=\frac{1}{24}$——这正是习题 2.9 第二条断言里的初始条件。]
% [Tutor Note: 同理，(2.64) 括号里的 $t_0^2/2$ 对应初始值 $\langle\tau_0^3\rangle_0=1$：设 $\langle\tau_0^3\rangle_0=A$，则 $Z=1+\hbar A t_0^3/6+\cdots$，$L_{-1}Z=0$ 在 $\hbar^2t_0^2$ 阶给出 $\hbar^2\cdot\frac{A}{2}-\hbar^2\cdot\frac{1}{2}=0$，故 $A=1$。一般 $t_0$ 次幂上得到的是弦方程（删去一个 $\tau_0$ 相当于把一个指标降 1，见 (2.57)）。]

算子 $L_{-1}$ 与 $L_0$ 可以看作（Virasoro 代数的一个子代数的表示的）开端。更确切地说，考虑由
\begin{equation}
\mathcal{L}_n=-z^{n+1}\frac{\partial}{\partial z}\,,\qquad n\ge-1. \tag{2.66}
\end{equation}
张成的全纯微分算子构成的 Lie 代数 $\mathrm{Vir}_{\ge-1}$。
% [Tutor Note: 记号辨析——原文用了三套 “L”：(i) 花体 $\mathcal{L}_n$，即 (2.66) 的抽象 Virasoro 生成元，其交换子无 $\hbar$：$[\mathcal{L}_m,\mathcal{L}_n]=(m-n)\mathcal{L}_{m+n}$；(ii) 斜体 $L_n$，(2.64)(2.65)(2.67) 的具体微分算子，定理 2.6 与 (2.69) 用的也是斜体；(iii) 正体 $\mathrm{L}_n$，习题 2.10 中重标度后的算子，满足带 $\hbar^2$ 的 (2.68)。译文按原文字形保留，勿混用。]

% [原文 p.22]
其括号由 $[\mathcal{L}_m,\mathcal{L}_n]=(m-n)\mathcal{L}_{m+n}$ 给出。

微分算子构成的集合 $(L_{-1},L_0)$ 可以唯一地（在一定的齐次性限制下）延拓为 $\mathrm{Vir}_{\ge-1}$（的一个 $\hbar$-形变）的完全表示。对 $n\ge1$，它们由下式给出
\begin{align}
L_n &= \hbar\frac{\partial}{\partial t_{n+1}}-\hbar^2\Bigg(\sum_{k\ge0}\frac{(2n+2k+1)!!}{(2n+3)!!(2k-1)!!}\,t_k\frac{\partial}{\partial t_{k+n}} \nonumber\\
&\qquad+\frac{1}{2}\sum_{\substack{a,b\ge0\\ a+b=n-1}}\frac{(2a+1)!!(2b+1)!!}{(2n+3)!!}\,\frac{\partial^2}{\partial t_a\partial t_b}\Bigg). \tag{2.67}
\end{align}
这里 $m!!$ 表示双阶乘，按 $m!!=m\cdot(m-2)!!$ 递归定义，初始条件为 $0!!=1!!=1$。这些正是在 Bouchard 的讲义 [Bou26] 中出现的微分约束！！
% [Tutor Note: 原文的两处细节照录：一是 “Bouchard 的讲义 [Bou26]！！” 结尾的双感叹号（源文件遗留或作者强调），二是 (2.67) 明确只对 $n\ge1$ 成立。核对 $n=0$：(2.67) 第一项因 $\frac{(2k+1)!!}{3!!(2k-1)!!}=\frac{2k+1}{3}$ 与 (2.65) 一致，但第二个求和为空、且缺常数 $+\frac{1}{24}$，所以 (2.65) 不能由 (2.67) 令 $n=0$ 得到。]
% [Tutor Note: 核对 (2.67) 的第二项在 $n=1$ 时：$a+b=n-1=0$ 只有 $(a,b)=(0,0)$，系数为 $\frac{1}{2}\cdot\frac{1!!\cdot1!!}{5!!}=\frac{1}{30}$，即 $L_1$ 含 $\frac{1}{30}\partial_0^2$。这个 $1/30$ 与 (2.70) 中二次项的系数 $\frac{1}{2}\frac{(2a+1)!!(2b+1)!!}{(2d_1+1)!!}$ 在 $d_1=2,a=b=0$ 时相同（$\frac{1}{2}\cdot\frac{1}{5!!}=\frac{1}{30}$），两式互相印证。]

\begin{exercise}[2.10]
证明由方程 (2.64)、(2.65) 与 (2.67) 定义的微分算子集合 $\bigl(\mathrm{L}_n:=-\frac{(2n+3)!!}{2}L_n\bigr)_{n\ge-1}$ 确实构成 $\mathrm{Vir}_{\ge-1}$ 的表示：
\begin{equation}
[\mathrm{L}_m,\mathrm{L}_n]=\hbar^2(m-n)\mathrm{L}_{m+n}. \tag{2.68}
\end{equation}
这一点连同算子的形式 (2.67)，证明了 $(\mathrm{L}_n)_{n\ge-1}$ 构成一个 Airy 理想 [Bou26]。
\end{exercise}
% [Tutor Note: 用 $(m,n)=(0,-1)$ 可显式核对 (2.68)。把 (2.64)(2.65) 直接作用在 $t_0$ 上算得 $[L_0,L_{-1}]=-\frac{2}{3}\hbar^2L_{-1}$；再乘上 $c_0c_{-1}=\bigl(-\frac{3!!}{2}\bigr)\bigl(-\frac{1!!}{2}\bigr)=\frac{3}{4}$，得 $[\mathrm{L}_0,\mathrm{L}_{-1}]=-\frac{1}{2}\hbar^2L_{-1}=\hbar^2\mathrm{L}_{-1}$，正是 (2.68) 在 $(0,-1)$ 处的形式（注意 $\mathrm{L}_{-1}=-\frac{1!!}{2}L_{-1}$）。这说明 $c_n=-\frac{(2n+3)!!}{2}$ 的作用是把 (2.67) 中 $(2n+3)!!$ 的分母吸收掉，使交换子重新闭合。]

\begin{theorem}[2.6]
（Witten 猜想/Kontsevich 定理）微分算子 $(L_n)_{n\ge-1}$ 湮灭配分函数 $Z$：
\begin{equation}
L_nZ=0\qquad \forall n\ge-1. \tag{2.69}
\end{equation}
此外，上述方程组（称为 Virasoro 约束）唯一地确定所有相交数。
\end{theorem}
% [Tutor Note: 定理中的斜体 $L_n$ 与习题 2.10 的正体 $\mathrm{L}_n$ 只差非零常数 $-\frac{(2n+3)!!}{2}$，因此 “$L_nZ=0$ 对所有 $n$” 与 “$\mathrm{L}_nZ=0$ 对所有 $n$” 完全等价；原文在定理里回到斜体记号，是因为湮灭性质与标度无关。历史顺序上，Witten 在 [Wit91] 提出猜想，Kontsevich 在 [Kon92] 用矩阵模型证明；(2.69) 这种 Virasoro 形式与 KdV 表述的等价性由 [DVV91] 证明，见下一段。]

我们指出，Witten 对其猜想的原始表述是：$Z$ 是 Korteweg--de Vries（KdV）层级的唯一一个满足弦方程 $L_{-1}Z=0$ 的 tau 函数。KdV 层级是一个偏微分方程的无穷序列，它在某种意义上推广了 KdV 方程。用 Virasoro 约束表述的等价命题由 R. Dijkgraaf、H. Verlinde、E. Verlinde [DVV91] 证明。

\begin{exercise}[2.11]
证明 Virasoro 约束等价于 Witten 关联函数的如下拓扑递归：
\begin{align}
\langle\tau_{d_1}\cdots\tau_{d_n}\rangle_g
&= \sum_{m=2}^{n}\frac{(2d_1+2d_m-1)!!}{(2d_1+1)!!\,(2d_m-1)!!}\,
\langle\tau_{d_1+d_m-1}\tau_{d_2}\cdots\widehat{\tau_{d_m}}\cdots\tau_{d_n}\rangle_g \nonumber\\
&\quad+\frac{1}{2}\sum_{a+b=d_1-2}\frac{(2a+1)!!\,(2b+1)!!}{(2d_1+1)!!}
\Bigl(\langle\tau_a\tau_b\tau_{d_2}\cdots\tau_{d_n}\rangle_{g-1} \nonumber\\
&\qquad+\sum_{\substack{g_1+g_2=g\\ I_1\sqcup I_2=\{d_2,\ldots,d_n\}}}
\langle\tau_a\tau_{I_1}\rangle_{g_1}\langle\tau_b\tau_{I_2}\rangle_{g_2}\Bigr). \tag{2.70}
\end{align}
% [Tutor Note: (2.70) 对第一个标记不对称：只有 $d_1\ge2$ 时第二个求和才非空，故实用中须取某个 $d_i\ge2$ 的标记充当 $d_1$。$g\ge2$ 时这样的标记必存在——因为 $\sum_i d_i=3g-3+n>n$，若所有 $d_i\le1$ 则 $\sum_i d_i\le n$，矛盾；而 $g\le1$ 的基始情形见下一条批注。按 $2g-2+n$ 归纳可证递归终止：第一项使 $n$ 减 1，第二项的第一种情形使 $g$ 减 1，第二种情形把标记分成两组，都严格减小 $2g-2+n$。]

% [原文 p.23]
\begin{table}[htbp]\centering\normalfont\footnotesize
\setlength{\tabcolsep}{4pt}
\begin{tabular}[t]{c l c}
\toprule
$(g,n)$ & $\langle\tau_{d_1}\cdots\tau_{d_n}\rangle_g$ & $*$\\
\midrule
$(0,3)$ & $\langle\tau_0^3\rangle_0$ & $1$\\
\midrule
$(0,4)$ & $\langle\tau_0^3\tau_1\rangle_0$ & $1$\\
\midrule
\multirow{2}{*}{$(0,5)$} & $\langle\tau_0^4\tau_2\rangle_0$ & $1$\\
 & $\langle\tau_0^3\tau_1^2\rangle_0$ & $2$\\
\midrule
\multirow{3}{*}{$(0,6)$} & $\langle\tau_0^5\tau_3\rangle_0$ & $1$\\
 & $\langle\tau_0^4\tau_1\tau_2\rangle_0$ & $3$\\
 & $\langle\tau_0^3\tau_1^3\rangle_0$ & $6$\\
\midrule
\multirow{5}{*}{$(0,7)$} & $\langle\tau_0^6\tau_4\rangle_0$ & $1$\\
 & $\langle\tau_0^5\tau_1\tau_3\rangle_0$ & $4$\\
 & $\langle\tau_0^5\tau_2^2\rangle_0$ & $6$\\
 & $\langle\tau_0^4\tau_1^2\tau_2\rangle_0$ & $12$\\
 & $\langle\tau_0^3\tau_1^4\rangle_0$ & $24$\\
\bottomrule
\end{tabular}\hspace{1.2em}%
\begin{tabular}[t]{c l c}
\toprule
$(g,n)$ & $\langle\tau_{d_1}\cdots\tau_{d_n}\rangle_g$ & $*$\\
\midrule
$(1,1)$ & $\langle\tau_1\rangle_1$ & $\frac{1}{24}$\\
\midrule
\multirow{2}{*}{$(1,2)$} & $\langle\tau_0\tau_2\rangle_1$ & $\frac{1}{24}$\\
 & $\langle\tau_1^2\rangle_1$ & $\frac{1}{24}$\\
\midrule
\multirow{3}{*}{$(1,3)$} & $\langle\tau_0^2\tau_3\rangle_1$ & $\frac{1}{24}$\\
 & $\langle\tau_0\tau_1\tau_2\rangle_1$ & $\frac{1}{12}$\\
 & $\langle\tau_1^3\rangle_1$ & $\frac{1}{12}$\\
\midrule
\multirow{5}{*}{$(1,4)$} & $\langle\tau_0^3\tau_4\rangle_1$ & $\frac{1}{24}$\\
 & $\langle\tau_0^2\tau_1\tau_3\rangle_1$ & $\frac{1}{8}$\\
 & $\langle\tau_0^2\tau_2^2\rangle_1$ & $\frac{1}{6}$\\
 & $\langle\tau_0\tau_1^2\tau_2\rangle_1$ & $\frac{1}{4}$\\
 & $\langle\tau_1^4\rangle_1$ & $\frac{1}{4}$\\
\bottomrule
\end{tabular}\hspace{1.2em}%
\begin{tabular}[t]{c l c}
\toprule
$(g,n)$ & $\langle\tau_{d_1}\cdots\tau_{d_n}\rangle_g$ & $*$\\
\midrule
$(2,1)$ & $\langle\tau_4\rangle_2$ & $\frac{1}{1152}$\\
\midrule
\multirow{3}{*}{$(2,2)$} & $\langle\tau_0\tau_5\rangle_2$ & $\frac{1}{1152}$\\
 & $\langle\tau_1\tau_4\rangle_2$ & $\frac{1}{384}$\\
 & $\langle\tau_2\tau_3\rangle_2$ & $\frac{29}{5760}$\\
\midrule
$(3,1)$ & $\langle\tau_7\rangle_3$ & $\frac{1}{82944}$\\
\midrule
\multirow{5}{*}{$(3,2)$} & $\langle\tau_0\tau_8\rangle_3$ & $\frac{1}{82944}$\\
 & $\langle\tau_1\tau_7\rangle_3$ & $\frac{5}{82944}$\\
 & $\langle\tau_2\tau_6\rangle_3$ & $\frac{77}{414720}$\\
 & $\langle\tau_3\tau_5\rangle_3$ & $\frac{503}{1451520}$\\
 & $\langle\tau_4^2\rangle_3$ & $\frac{607}{1451520}$\\
\midrule
$(4,1)$ & $\langle\tau_{10}\rangle_4$ & $\frac{1}{7962624}$\\
\bottomrule
\end{tabular}
\caption*{表 1：部分 $\psi$ 类相交数，用拓扑递归关系 (2.70) 算出。}
\end{table}
% [Tutor Note: 表 1 全部 34 个数值已用 (2.70) 逐格重算核对（基始情形取 $g=0$ 的闭合公式 (2.60) 与 $g=1$ 的 (2.61)），34/34 与表一致；其中 $(3,2)$ 行的 $\langle\tau_4^2\rangle_3=\frac{607}{1451520}$、$(2,2)$ 行的 $\langle\tau_2\tau_3\rangle_2=\frac{29}{5760}$ 是最易抄错的两位大分母，可优先用它们抽查。]
% [Tutor Note: 还可用弦方程/膨胀方程交叉核对若干格：弦方程（把一个指标降 1）给出 $\langle\tau_0\tau_8\rangle_3=\langle\tau_7\rangle_3=\frac{1}{82944}$、$\langle\tau_0\tau_5\rangle_2=\langle\tau_4\rangle_2=\frac{1}{1152}$；膨胀方程给出 $\langle\tau_1^4\rangle_1=(2\cdot1-2+3)\langle\tau_1^3\rangle_1=3\cdot\frac{1}{12}=\frac{1}{4}$。$(0,n)$ 各行就是 (2.60) 的多重组合数，例如 $\langle\tau_0^3\tau_1^4\rangle_0=\frac{(7-3)!}{0!^3\,1!^4}=24$、$\langle\tau_0^5\tau_2^2\rangle_0=\frac{4!}{0!^5\,2!^2}=6$。]
% [Tutor Note: 排版说明——原文表头第三列的 “$*$” 是数值列的表头（原文即如此，非笔误）；三个子表并列且各自带 \toprule/\bottomrule。译文用三张 [t] 对齐的 tabular 并排复现，同一 $(g,n)$ 用 multirow 跨行。另：Table 1 在原文是浮动体，排在 p23 顶部、插在习题 2.11 的 (2.70) 与下面“证明上述递归等价于……”一段之间；译文把表放在同一位置（习题 2.11 内部）。]

证明上述递归等价于 Airy 谱曲线 $\bigl(\PP^1,\,x(z)=\frac{z^2}{2},\,y(z)=z,\,\omega_{0,2}(z_1,z_2)=\frac{\dd z_1\dd z_2}{(z_1-z_2)^2}\bigr)$ 上的 Eynard--Orantin 拓扑递归公式 [EO07]（见 [Bou26]）：
\begin{equation}
\omega_{g,n}(z_1,\ldots,z_n)=(-1)^n\sum_{\substack{d_1,\ldots,d_n\ge0\\ d_1+\cdots+d_n=3g-3+n}}\langle\tau_{d_1}\cdots\tau_{d_n}\rangle_g\prod_{i=1}^{n}\frac{(2d_i+1)!!}{z_i^{2d_i+2}}\,\dd z_i. \tag{2.71}
\end{equation}
\end{exercise}
% [Tutor Note: (2.71) 的 $(-1)^n$ 可与 (2.73) 对证：对 (2.73) 的 $\prod_i\frac{(2d_i-1)!!}{\lambda_i^{2d_i+1}}$ 逐项求导 $\frac{\dd}{\dd\lambda_i}$ 得 $-\frac{(2d_i+1)!!}{\lambda_i^{2d_i+2}}$，$n$ 个因子共出 $(-1)^n$，与 (2.71) 的 $(-1)^n$ 逐个吻合（并把 $\lambda_i$ 换成 $z_i$）。]

正如引言中提到的，Witten 提出上述猜想的动机根源于二维量子引力。在经典情形下，时空是一个曲面，而引力场是该曲面上的黎曼度规。为了量子化这样的一个理论，人们应当在所有可能曲面上所有可能的黎曼度规构成的空间上计算某个路径积分。固定拓扑曲面上的黎曼度规空间是无穷维的，而有两条可能的途径来为这样一个不适定的量赋予意义。
\begin{itemize}
\item 第一条途径是用小三角形逼近黎曼曲面。于是，对所有度规的积分被替换为对三角剖分的求和。这个组合问题可以求解，而 Virasoro 约束正是在关于曲面三角剖分计数的研究中出现的，这类计数问题具有矩阵模型表述。
\item 另一条途径是，先对所有共形等价的度规积分，从而算出配分函数。之后，剩余的积分在黎曼曲面的模空间上进行；更确切地说，人们必须计算形如 $\langle\tau_{d_1}\cdots\tau_{d_n}\rangle_g$ 的积分。
\end{itemize}

% [原文 p.24]
Witten 猜想断言，由这两条途径得到的配分函数相互重合，其依据是“引力理论唯一”这一物理预期。

Kontsevich 的证明遵循矩阵模型/离散化的想法（严格证明见 [Zvo]）。他首先考虑亏格 $g$、$n$ 个具有固定长度 $L_1,\ldots,L_n$ 的面的度量带状图的模空间，该空间带有自然的（辛）体积形式。把度量带状图解释为黎曼度规的离散化，Kontsevich 把这些体积精确地表示为 $\psi$ 类相交数
\begin{align}
V_{g,n}(L_1,\ldots,L_n) &= \int_{\Mc_{g,n}}\exp\left(\frac{1}{2}\sum_{i=1}^{n}L_i^2\psi_i\right) \nonumber\\
&= \sum_{\substack{d_1,\ldots,d_n\ge0\\ d_1+\cdots+d_n=3g-3+n}}\langle\tau_{d_1}\cdots\tau_{d_n}\rangle_g\prod_{i=1}^{n}\frac{L_i^{2d_i}}{2^{d_i}d_i!}. \tag{2.72}
\end{align}
注意 $V_{g,n}(L_1,\ldots,L_n)$ 是边界长度平方的对称多项式。这种体积的 Laplace 变换可算得为有理函数
\begin{align}
\widehat{V}_{g,n}(\lambda_1,\ldots,\lambda_n) &= \left(\prod_{i=1}^{n}\int_0^\infty \dd L_i\,\ee^{-\lambda_iL_i}\right)V_{g,n}(L_1,\ldots,L_n) \nonumber\\
&= \sum_{\substack{d_1,\ldots,d_n\ge0\\ d_1+\cdots+d_n=3g-3+n}}\langle\tau_{d_1}\cdots\tau_{d_n}\rangle_g\prod_{i=1}^{n}\frac{(2d_i-1)!!}{\lambda_i^{2d_i+1}}. \tag{2.73}
\end{align}
% [Tutor Note: 由 (2.72) 到 (2.73) 用的是 $\int_0^\infty\dd L\,\ee^{-\lambda L}L^{2d}=\frac{(2d)!}{\lambda^{2d+1}}$ 以及恒等式 $\frac{(2d)!}{2^d d!}=(2d-1)!!$；正是这个恒等式把 (2.72) 分母中的 $2^{d_i}d_i!$ 化成 (2.73) 分子中的 $(2d_i-1)!!$。取 $d=1$ 检验：$\frac{2!}{2\cdot1}=1=1!!$；取 $d=2$：$\frac{4!}{4\cdot2}=3=3!!$。]
注意 $(\dd_{\lambda_1}\cdots\dd_{\lambda_n})\widehat{V}_{g,n}(\lambda)=\omega_{g,n}(\lambda)$ 恰是由 Airy 谱曲线计算出的 (2.71) 中的拓扑递归关联函数。

由于 $V_{g,n}(L_1,\ldots,L_n)$ 是亏格 $g$、$n$ 个具有固定长度 $L_1,\ldots,L_n$ 的面的度量带状图模空间的体积，他把 Laplace 变换表示为对带状图的求和：
\begin{equation}
\widehat{V}_{g,n}(\lambda_1,\ldots,\lambda_n)=2^{2g-2+n}\sum_{\mathbb{G}}\frac{1}{|\Aut(\mathbb{G})|}\prod_{e=(i,j)\in E(\mathbb{G})}\frac{1}{\lambda_i+\lambda_j}, \tag{2.74}
\end{equation}
其中求和遍历亏格 $g$、$n$ 个面标号为 $1,\ldots,n$ 的所有三价带状图，而 $e=(i,j)$ 表示图中与边 $e$ 相邻的两个（可能相同的）面。

例如，取 $g=n=1$。此时只有一个三价带状图，见图 13，它的自同构群为 $\Z_6$（三条边的循环置换，以及交换两个顶点的额外 $\Z_2$ 对称）。
\figomit{13}{唯一的 $(1,1)$ 型三价带状图。}

% [原文 p.25]
于是 Kontsevich 公式 (2.74) 给出
\begin{equation}
\widehat{V}_{1,1}(\lambda_1)=2\cdot\frac{1}{6}\cdot\left(\frac{1}{2\lambda_1}\right)^{3}=\frac{1}{24}\frac{1}{\lambda_1^3}, \tag{2.75}
\end{equation}
这确实给出 $\langle\tau_1\rangle_1=\frac{1}{24}$，与 (2.73) 相符。
% [Tutor Note: 数值核对 (2.74)→(2.75)：$2^{2g-2+n}\big|_{g=n=1}=2$，$|\Aut(\mathbb{G})|=6$，三条边各配因子 $\frac{1}{\lambda_1+\lambda_1}=\frac{1}{2\lambda_1}$，故 $\widehat V_{1,1}=2\cdot\frac{1}{6}\cdot\frac{1}{(2\lambda_1)^3}=\frac{1}{24\lambda_1^3}$；另一方面 (2.73) 中 $d_1=1$ 的项为 $\frac{(2-1)!!}{\lambda_1^3}\langle\tau_1\rangle_1=\frac{1}{24\lambda_1^3}$，两者一致，从而再次得到 $\langle\tau_1\rangle_1=\frac{1}{24}$。]

一方面，Kontsevich 定理给出图的求和，其中每个图以其对称因子和边的权之积加权。这正是由 Wick 定理得到的典型图求和，因此它可以由对高斯 Hermite 矩阵积分的一个微扰得到。具体地说，三价带状图由三次形式的矩阵积分生成，即所谓的 \emph{Airy 矩阵积分}：
\begin{equation}
Z(\Lambda)=\frac{1}{Z_0(\Lambda)}\int \dd X\exp\left(N\,\mathrm{tr}\left[\frac{X^3}{3}-\Lambda X^2\right]\right),\qquad \Lambda=\mathrm{diag}(\lambda_1,\ldots,\lambda_N). \tag{2.76}
\end{equation}
这里 $Z_0(\Lambda)=(\pi/N)^{N^2/2}\prod_{i,j}(\lambda_i+\lambda_j)^{-1/2}$ 是一个归一化常数。由 Wick 定理，可以把 $\log Z(\Lambda)$ 的大 $N$ 展开写成对三价带状图的求和：
\begin{equation}
\log Z(\Lambda)=\sum_{\substack{g\ge0,\,n\ge1\\ 2g-2+n>0}}\frac{N^{-(2g-2+n)}}{n!}\sum_{\mathbb{G}}\frac{1}{|\Aut(\mathbb{G})|}\prod_{e=(i,j)\in E(\mathbb{G})}\frac{1}{\lambda_i+\lambda_j}, \tag{2.77}
\end{equation}
其中求和遍历亏格 $g$、$n$ 个标号面的所有三价带状图。
% [Tutor Note: $Z$ 的定义里已除掉 $Z_0(\Lambda)$，所以 (2.77) 的 $\log Z$ 没有 $n=0$ 的常数项，这正是其求和条件也写成 $2g-2+n>0$ 的原因。另注意原文的 $\mathrm{tr}$ 是小写正体（与 $\Tr$ 不同），(2.76) 与 (2.78) 一致用小写。]
% [Tutor Note: 核对 (2.77) 与 (2.63) 的相容性：把 $\hbar=2/N$ 代入 (2.63) 得系数 $N^{-(2g-2+n)}\,2^{2g-2+n}$，而 (2.74) 前恰好有因子 $2^{2g-2+n}$；两者相消后正是 (2.77) 的 $N^{-(2g-2+n)}$。这说明 (2.77) 与 (2.63)(2.74) 连同时间变量的识别 $t_d=\mathrm{tr}(\Lambda^{-2d-1})/(2d-1)!!$ 是自洽的。]

最后，分部积分（在此语境下也称为 \emph{Schwinger--Dyson 方程}）表明，在取 $\hbar=2/N$、并把时间变量取为矩阵模型展开中自然出现的 $\Lambda$ 的归一化迹之后，$Z(\Lambda)$ 满足 Virasoro 约束 (2.69)：
\begin{equation}
t_d=\frac{\mathrm{tr}(\Lambda^{-2d-1})}{(2d-1)!!},\qquad d\ge0. \tag{2.78}
\end{equation}
% [Tutor Note: 取 $d=0$ 核对 (2.78)：因 $(-1)!!=1$，得 $t_0=\mathrm{tr}(\Lambda^{-1})$，与 (2.73) 中 $d_i=0$ 的因子 $\frac{(-1)!!}{\lambda_i}=\frac{1}{\lambda_i}$ 相合（把 $\frac{1}{\lambda_i}$ 换成 $\mathrm{tr}$ 的求和项即得 $t_0$）。]
\figomit{14}{按领头渐近归一化的两点关联函数：注意它们收敛到 $1+\mathrm{O}(g^{-1})$。同时注意带 $\tau_0$ 插入（绿色）与不带 $\tau_0$ 插入（蓝色）的关联函数具有不同的收敛行为；这提示次领头项确实依赖于分拆 $(d_1,\ldots,d_n)$。事实确实如此，并且可以借助复兴理论证明。}

% [原文 p.26]
值得指出的是，通过复兴（resurgence）技术（见 [ABS19] 或 I. Aniceto 与 M. Mariño 的讲义 [Ani; Mar26]），可以计算大亏格渐近的 Witten 关联函数 [EGGGL]，见图 14：
\begin{equation}
\langle\tau_{d_1}\cdots\tau_{d_n}\rangle_g\prod_{i=1}^{n}(2d_i+1)!!=\frac{2^{n-1}}{2\pi}\,\frac{\Gamma(2g-2+n)}{\left(\frac{2}{3}\right)^{2g-2+n}}\left(1+\mathrm{O}(g^{-1})\right). \tag{2.79}
\end{equation}
% [Tutor Note: 用表 1 的末行核对 (2.79)：$g=4,n=1,d_1=10$ 时左端为 $\langle\tau_{10}\rangle_4\cdot21!!=\frac{13749310575}{7962624}\approx1726.73$，右端为 $\frac{1}{2\pi}\Gamma(7)\bigl(\frac{3}{2}\bigr)^{7}\approx1957.90$，比值约 $0.882$；$g=4$ 还太小，不足以看出 $1+\mathrm{O}(g^{-1})$。改用闭式 $\langle\tau_{3g-2}\rangle_g=\frac{1}{24^g g!}$ 逐 $g$ 比较，比值依次为 $0.524,\ 0.764,\ 0.843,\ 0.882,\ 0.906,\ 0.921,\ \ldots,\ 0.961\ (g=12)$，单调趋 1，与 (2.79) 相容，也解释了图 14（从略）中纵轴落在 $10^0$ 附近的原因。]
次领头渐近也可以借助复兴得到。这一结果的第一个证明使用了组合与概率论证，归功于 A. Aggarwal [Agg21]。注意 Stokes 常数 $S=\ii$ 与瞬子作用量 $A=2/3$，它们与 Airy 函数的相应量一致。这自然不是巧合！
